\section{Off-Policy 策略梯度}

\subsection{On-Policy 的局限}

Vanilla policy gradient 是完全 on-policy 的：每次梯度更新后，之前的数据就不能再用了。这意味着\textbf{每个梯度步都需要重新收集数据}，效率很低。

\subsection{Importance Sampling}

可以用 importance weights 来复用旧策略 $\pi_\theta$ 收集的数据来更新新策略 $\pi_{\theta'}$：
$$\nabla_{\theta'} J(\theta') \approx \frac{1}{N}\sum_{i=1}^{N}\sum_{t=1}^{T} \frac{\pi_{\theta'}(a_{i,t} \mid s_{i,t})}{\pi_\theta(a_{i,t} \mid s_{i,t})} \nabla_{\theta'} \log \pi_{\theta'}(a_{i,t} \mid s_{i,t}) \left(\sum_{t'=t}^{T} r(s_{i,t'}, a_{i,t'}) - b\right)$$

\begin{warningbox}{Importance Sampling 不是万能的}
当新旧策略差异很大时，importance weights 的方差会非常大，使得估计不可靠。因此 importance sampling 只适用于策略\textbf{非常接近}的情况。
\end{warningbox}

\subsection{Replay Buffer 与 Truncated Importance Sampling}
为了进一步提高数据利用率，可以把过去的 trajectory 存入 replay buffer，结合 truncated importance sampling（例如 PPO 的 clips / ratio clipping）来约束 weight：
$$\hat{g} = \min\left(\frac{\pi_{\theta'}(a\mid s)}{\pi_\theta(a\mid s)}, 1+\epsilon\right) \nabla_\theta \log \pi_\theta(a\mid s) \hat{A}.$$
Clip 机制限制了 ratio 的过度增长，防止过大的 importance weight 导致梯度爆炸。

\begin{importantbox}{Clip ratio 的作用}
\begin{itemize}
    \item 限制 ratio 在 $[1-\epsilon,1+\epsilon]$ 之间，避免 over-optimistic update；
    \item 仍保留了 off-policy 的复用能力，只要新旧策略保持相似。
\end{itemize}
\end{importantbox}

\subsection{方差控制提醒}
除了 ratio clipping，也可以 anneal $\epsilon$、增加 entropy bonus、加入 trust region gate。更复杂的做法是把每条 trajectory 关联的 importance weight 归一化，使得多样本 update 时不会被某一条高 weight 数据主导。

\subsection{本章小结}

Vanilla policy gradient 是 on-policy 的，数据利用率低。Importance sampling 可以实现一定程度的 off-policy 更新，但受限于新旧策略的相似度。

